\Chapter{28}

\Verse{1}
{
    \zA{话说林黛玉只因昨夜晴雯不开门一事，错疑在宝玉身上。 次日又可巧遇见饯花 之期，正在一腔无明未曾发泄，又勾起伤春愁思，因把些残花落瓣去掩埋，由不得
        感花伤己，哭了几声，便随口念了几句。 不想宝玉在山坡上听见，先不过点头感叹； 次又听到“侬今葬花人笑痴，他年葬侬知是谁”、“一朝春尽红颜老，花落人亡两
        不知”等句，不觉恸倒山坡上，怀里兜的落花撒了一地。 试想林黛玉的花颜月貌， 将来亦到无可寻觅之时，宁不心碎肠断？}
}
{
    \eA{Lin Tai-yü, the story goes, dwelt, after Ch'ing Wen's refusal, the previous
        night, to open the door, under the impression that the blame lay with
        Pao-yü.}
}
{
}

\Verse{2}
{
    \zA{既黛玉终归无可寻觅之时，推之于他人， 如宝钗、香菱、袭人等，亦可以到无可寻觅之时矣。 宝钗等终归无可寻觅之时，则 自己又安在呢？
        且自身尚不知何在何往，将来斯处、斯园、斯花、斯柳，又不知当 属谁姓？ 因此一而二二而三反复推求了去，真不知此时此际如何解释这段悲伤!正 是：
        花影不离身左右，鸟声只在耳东西。 那黛玉正自伤感，忽听山坡上也有悲声，心下想道：“人人都笑我有痴病，难 道还有一个痴的不成？”}
}
{
    \eA{The following day, which by another remarkable coincidence, happened to
        correspond with the season, when the god of flowers had to be feasted, her
        total ignorance of the true circumstances, and her resentment, as yet
        unspent, aroused again in her despondent thoughts, suggested by the decline
        of spring time. She consequently gathered a quantity of faded flowers and
        fallen petals, and went and interred them.}
}
{
}

\Verse{3}
{
    \zA{抬头一看，见是宝玉，黛玉便啐道：“呸!我打量是谁， 原来是这个狠心短命的――”刚说到“短命”二字，又把口掩住，长叹一声，自己 抽身便走。
        这里宝玉悲恸了一回，见黛玉去了，便知黛玉看见他躲开了，自己也觉无味。 抖抖土起来，下山寻归旧路，往怡红院来。 可巧看见黛玉在前头走，连忙赶上去，
        说道：“你且站着。 我知道你不理我； 我只说一句话，从今以后撩开手。” 黛玉回 头见是宝玉，待要不理他，听他说只说一句话，便道：“请说。”}
}
{
    \eA{Unable to check the emotion, caused by the decay of the flowers, she
        spontaneously recited, after giving way to several loud lamentations, those
        verses which Pao-yü, she little thought, overheard from his position on the
        mound. At first, he did no more than nod his head and heave sighs, full of
        feeling.}
}
{
}

\Verse{4}
{
    \zA{宝玉笑道：“两 句话，说了你听不听呢？” 黛玉听说，回头就走。 宝玉在身后面叹道：“既有今日， 何必当初？”
        黛玉听见这话，由不得站住，回头道：“当初怎么样？ 今日怎么样？” 宝玉道：“嗳!当初姑娘来了，那不是我陪着玩笑？ 凭我心爱的，姑娘要就拿去； 我
        爱吃的，听见姑娘也爱吃，连忙收拾的干干净净收着，等着姑娘回来。 一个桌子上 吃饭，一个床儿上睡觉。 丫头们想不到的，我怕姑娘生气，替丫头们都想到了。}
}
{
    \eA{But when subsequently his ear caught:  "Here I am fain these flowers to
        inter, but humankind will laugh me  as  a fool; Who knows who will, in years
        to come, commit me to my grave! In a twinkle springtime draws to an end, and
        maidens wax in age. Flowers fade and maidens die; and of either naught any
        more is  known."}
}
{
}

\Verse{5}
{
    \zA{我 想着姊妹们从小儿长大，亲也罢，热也罢，和气到了儿，才见得比别人好。 如今谁 承望姑娘人大心大，不把我放在眼里，三日不理、四日不见的，倒把外四路儿的什
        么‘宝姐姐’‘凤姐姐’的放在心坎儿上。 我又没个亲兄弟、亲妹妹，虽然有两个， 你难道不知道是我隔母的？ 我也和你是独出，只怕你和我的心一样。
        谁知我是白操 了这一番心，有冤无处诉！” 说着，不觉哭起来。}
}
{
    \eA{he unconsciously was so overpowered with grief that he threw himself on the
        mound, bestrewing the whole ground with the fallen flowers he carried in his
        coat, close to his chest. "When Tai-yü's flowerlike charms and moon-like
        beauty," he reflected, "by and bye likewise reach a time when they will
        vanish beyond any hope of recovery, won't my heart be lacerated and my
        feelings be mangled!}
}
{
}

\Verse{6}
{
    \zA{那时黛玉耳内听了这话，眼内见了这光景，心内不觉灰了大半，也不觉滴下泪 来，低头不语。 宝玉见这般形象，遂又说道：“我也知道我如今不好了，但只任凭
        我怎么不好，万不敢在妹妹跟前有错处。 就有一二分错处，你或是教导我，戒我下 次，或骂我几句，打我几下，我都不灰心。 谁知你总不理我，叫我摸不着头脑儿，
        少魂失魄，不知怎么样才好。 就是死了也是个屈死鬼，任凭高僧高道忏悔，也不能 超生，还得你说明了原故，我才得托生呢！”}
}
{
    \eA{And extending, since Tai-yü must at length some day revert to a state when
        it will be difficult to find her, this reasoning to other persons, like
        Pao-ch'ai, Hsiang Ling, Hsi Jen and the other girls, they too are equally
        liable to attain a state beyond the reach of human search. But when
        Pao-ch'ai and all the rest have ultimately reached that stage when no trace
        will be visible of them, where shall I myself be then?}
}
{
}

\Verse{7}
{
    \zA{黛玉听了这话，不觉将昨晚的事都忘在九霄云外了，便说道：“你既这么说， 为什么我去了，你不叫丫头开门呢！” 宝玉诧异道：“这话从那里说起？ 我要是这
        么着，立刻就死了！” 黛玉啐道：“大清早起‘死’呀‘活’的，也不忌讳!你说 有呢就有，没有就没有，起什么誓呢！” 宝玉道：“实在没有见你去，就是宝姐姐
        坐了一坐，就出来了。” 黛玉想了一想，笑道：“是了：必是丫头们懒怠动，丧声 歪气的，也是有的。” 宝玉道：“想必是这个原故。}
}
{
    \eA{And when my own human form will have vanished and gone, whither I know not
        yet, to what person, I wonder, will this place, this garden and these
        plants, revert?" From one to a second, and from a second to a third, he thus
        pursued his reflections, backwards and forwards, until he really did not
        know how he could best, at this time and at such a juncture, dispel his fit
        of anguish.}
}
{
}

\Verse{8}
{
    \zA{等我回去问了是谁，教训教训 他们就好了。” 黛玉道：“你的那些姑娘们，也该教训教训。 只是论理我不该说。
        今儿得罪了我的事小，倘或明儿‘宝姑娘’来，什么‘贝姑娘’来，也得罪了，事 情可就大了。” 说着抿着嘴儿笑。 宝玉听了，又是咬牙，又是笑。
        二人正说话，见丫头来请吃饭，遂都往前头来了。 王夫人见了黛玉，因问道： “大姑娘，你吃那鲍太医的药可好些？” 黛玉道：“也不过这么着。 老太太还叫我
        吃王大夫的药呢。”}
}
{
    \eA{His state is adequately described by:  The shadow of a flower cannot err
        from the flower itself to the left  or the right. The song of birds can only
        penetrate into the ear from the east or  the  west. Lin Tai-yü was herself a
        prey to emotion and agitation, when unawares sorrowful accents also struck
        her ear, from the direction of the mound.}
}
{
}

\Verse{9}
{
    \zA{宝玉道：“太太不知道：林妹妹是内症，先天生的弱，所以禁 不住一点儿风寒； 不过吃两剂煎药，疏散了风寒，还是吃丸药的好。” 王夫人道：
        “前儿大夫说了个丸药的名字，我也忘了。” 宝玉道：“我知道那些丸药，不过叫 他吃什么人参养荣丸。” 王夫人道：“不是。” 宝玉又道：“八珍益母丸？
        左归， 右归？ 再不就是八味地黄丸？” 王夫人道：“都不是。 我只记得有个‘金刚’两个 字的。”}
}
{
    \eA{"Every one," she cogitated, "laughs at me for labouring under a foolish
        mania, but is there likely another fool besides myself?" She then raised her
        head, and, casting a glance about her, she discovered that it was Pao-yü.
        "Ts'ui!" eagerly cried Tai-yü, "I was wondering who it was; but is it truly
        this ruthless-hearted and short-lived fellow!"}
}
{
}

\Verse{10}
{
    \zA{宝玉拍手笑道：“从来没听见有个什么‘金刚丸’!若有了‘金刚丸’， 自然有‘菩萨散’了！” 说的满屋里人都笑了。 宝钗抿嘴笑道：“想是天王补心丹。”
        王夫人笑道：“是这个名儿。 如今我也糊涂了。” 宝玉道：“太太倒不糊涂，都是 叫‘金刚’‘菩萨’支使糊涂了。”
        王夫人道：“扯你娘的臊!又欠你老子捶你了。” 宝玉笑道：“我老子再不为这个捶我。” 王夫人又道：“既有这个名儿，明儿就叫人买些来吃。”
        宝玉道：“这些药都 是不中用的。}
}
{
    \eA{But the moment the two words "short-lived" dropped from her mouth, she
        sealed her lips; and, heaving a deep sigh, she turned herself round and
        hurriedly walked off. Pao-yü, meanwhile, remained for a time a prey to
        melancholy. But perceiving that Tai-yü had retired, he at once realised that
        she must have caught sight of him and got out of his way;}
}
{
}

\Verse{11}
{
    \zA{太太给我三百六十两银子，我替妹妹配一料丸药，包管一料不完就好 了。” 王夫人道：“放屁!什么药就这么贵？” 宝玉笑道：“当真的呢。 我这个方
        子比别的不同，那个药名儿也古怪，一时也说不清，只讲那头胎紫河车，人形带叶 参，三百六十两不足。 龟大何首乌，千年松根茯苓胆，诸如此类的药不算为奇，只
        在群药里算。 那为君的药，说起来，唬人一跳!前年薛大哥哥求了我一二年，我才 给了他这方子。 他拿了方子去，又寻了二三年，花了有上千的银子，才配成了。}
}
{
    \eA{and, as his own company afforded him no pleasure, he shook the dust off his
        clothes, rose to his feet and descending the hill, he started for the I Hung
        court by the path by which he had come. But he espied Tai-yü walking in
        advance of him, and with rapid stride, he overtook her. "Stop a little!" he
        cried. "I know you don't care a rap for me;}
}
{
}

\Verse{12}
{
    \zA{太 太不信，只问宝姐姐。” 宝钗听说，笑着摇手儿说道：“我不知道，也没听见。 你 别叫姨娘问我。” 王夫人笑道：“到底是宝丫头好孩子，不撒谎。”
        宝玉站在当地， 听见如此说，一回身把手一拍，说道：“我说的倒是真话呢，倒说撒谎！” 口里说
        着，忽一回身，只见林黛玉坐在宝钗身后抿着嘴笑，用手指头在脸上画着羞他。 凤姐因在里间屋里看着人放桌子，听如此说，便走来笑道：“宝兄弟不是撒谎，
        这倒是有的。}
}
{
    \eA{but I'll just make one single remark, and from this day forward we'll part
        company." Tai-yü looked round. Observing that it was Pao-yü, she was about
        to ignore him; hearing him however mention that he had only one thing to
        say, "Please tell me what it is," she forthwith rejoined. Pao-yü smiled at
        her. "If I pass two remarks will you listen to me; yes or no?" he asked.}
}
{
}

\Verse{13}
{
    \zA{前日薛大爷亲自和我来寻珍珠，我问他做什么，他说配药。 他还抱怨 说：‘不配也罢了，如今那里知道这么费事！’ 我问：‘什么药？’ 他说是宝兄弟
        说的方子，说了多少药，我也不记得。 他又说：‘不是我就买几颗珍珠了，只是必 要头上戴过的，所以才来寻几颗。 要没有散的花儿，就是头上戴过的拆下来也使得。
        过后儿我拣好的再给穿了来。’ 我没法儿，只得把两枝珠子花儿现拆了给他。 还要 一块三尺长、上用的大红纱，拿乳钵研了面子呢。”}
}
{
    \eA{At these words, Tai-yü twisted herself round and beat a retreat. Pao-yü
        however followed behind. "Since this is what we've come to now," he sighed,
        "what was the use of what existed between us in days gone by?" As soon as
        Tai-yü heard his exclamation, she stopped short impulsively. Turning her
        face towards him, "what about days gone by," she remarked, "and what about
        now?" "Ai!"}
}
{
}

\Verse{14}
{
    \zA{凤姐说一句，宝玉念一句佛。 凤姐说完了，宝玉又道：“太太打量怎么着？ 这不过也是将就罢咧。 正经按方子，
        这珍珠宝石是要在古坟里找，有那古时富贵人家儿装裹的头面拿了来才好。 如今那 里为这个去刨坟掘墓？ 所以只是活人带过的也使得。”
        王夫人听了道：“阿弥陀佛， 不当家花拉的!就是坟里有，人家死了几百年，这会子翻尸倒骨的，作了药也不灵 啊。” 宝玉因向黛玉道：“你听见了没有？
        难道二姐姐也跟着我撒谎不成？”}
}
{
    \eA{ejaculated Pao-yü, "when you got here in days gone by, wasn't I your
        playmate in all your romps and in all your fun? My heart may have been set
        upon anything, but if you wanted it you could take it away at once. I may
        have been fond of any eatable, but if I came to learn that you too fancied
        it, I there and then put away what could be put away, in a clean place, to
        wait, Miss, for your return.}
}
{
}

\Verse{15}
{
    \zA{脸望着 黛玉说，却拿眼睛瞟着宝钗。 黛玉便拉王夫人道：“舅母听听，宝姐姐不替他圆谎， 他只问着我！” 王夫人也道：“宝玉很会欺负你妹妹。”
        宝玉笑道：“太太不知道 这个原故。 宝姐姐先在家里住着，薛大哥的事他也不知道，何况如今在里头住着呢？ 自然是越发不知道了。
        林妹妹才在背后，以为是我撒谎，就羞我。” 正说着，见贾母房里的丫头找宝玉和黛玉去吃饭。 黛玉也不叫宝玉，便起身带 着那丫头走。
        那丫头说：“等着宝二爷一块儿走啊。”}
}
{
    \eA{We had our meals at one table; we slept in one and the same bed; whatever
        the servant-girls could not remember, I reminded them of, for fear lest your
        temper, Miss, should get ruffled.}
}
{
}

\Verse{16}
{
    \zA{黛玉道：“他不吃饭，不和 咱们走，我先走了。” 说着，便出去了。 宝玉道：“我今儿还跟着太太吃罢。” 王
        夫人道：“罢罢，我今儿吃斋，你正经吃你的去罢。” 宝玉道：“我也跟着吃斋。” 说着，便叫那丫头：“去罢。” 自己跑到桌子上坐了。
        王夫人向宝钗等笑道：“你 们只管吃你们的，由他去罢。” 宝钗因笑道：“你正经去罢。 吃不吃，陪着林妹妹 走一趟，他心里正不自在呢。 何苦来？”
        宝玉道：“理他呢，过一会子就好了。”}
}
{
    \eA{I flattered myself that cousins, who have grown up together from their
        infancy, as you and I have, would have continued, through intimacy or
        friendship, either would have done, in peace and harmony until the end, so
        as to make it palpable that we are above the rest. But, contrary to all my
        expectations, now that you, Miss, have developed in body as well as in mind,
        you don't take the least heed of me.}
}
{
}

\Verse{17}
{
    \zA{一时吃过饭，宝玉一则怕贾母惦记，二则也想着黛玉，忙忙的要茶漱口。 探春 惜春都笑道：“二哥哥，你成日家忙的是什么？ 吃饭吃茶也是这么忙碌碌的。” 宝
        钗笑道：“你叫他快吃了瞧黛玉妹妹去罢。 叫他在这里胡闹什么呢？” 宝玉吃了茶 便出来，一直往西院来。 可巧走到凤姐儿院前，只见凤姐儿在门前站着，蹬着门槛
        子，拿耳挖子剔牙，看着十来个小厮们挪花盆呢。 见宝玉来了，笑道：“你来的好， 进来，进来，替我写几个字儿。” 宝玉只得跟了进来。}
}
{
    \eA{You lay hold instead of some cousin Pao or cousin Feng or other from here,
        there and everywhere and give them a place in your affections; while on the
        contrary you disregard me for three days at a stretch and decline to see
        anything of me for four! I have besides no brother or sister of the same
        mother as myself. It's true there are a couple of them, but these, are you
        not forsooth aware, are by another mother!}
}
{
}

\Verse{18}
{
    \zA{到了房里，凤姐命人取过笔 砚纸来，向宝玉道：“大红妆缎四十匹，蟒缎四十匹，各色上用纱一百匹，金项圈 四个。” 宝玉道：“这算什么？
        又不是帐，又不是礼物，怎么个写法儿？” 凤姐儿 道：“你只管写上，横竖我自己明白就罢了。” 宝玉听说，只得写了。 凤姐一面收
        起来，一面笑道：“还有句话告诉你，不知依不依？ 你屋里有个丫头叫小红的，我 要叫了来使唤，明儿我再替你挑一个，可使得么？”}
}
{
    \eA{You and I are only children, so I ventured to hope that you would have
        reciprocated my feelings. But, who'd have thought it, I've simply thrown
        away this heart of mine, and here I am with plenty of woes to bear, but with
        nowhere to go and utter them!" While expressing these sentiments, tears,
        unexpectedly, trickled from his eyes.}
}
{
}

\Verse{19}
{
    \zA{宝玉道：“我屋里的人也多的 很，姐姐喜欢谁，只管叫了来，何必问我？” 凤姐笑道：“既这么着，我就叫人带 他去了。” 宝玉道：“只管带去罢。” 说着要走。
        凤姐道：“你回来，我还有一句 话呢。” 宝玉道：“老太太叫我呢，有话等回来罢。” 说着，便至贾母这边。 只见都已吃完了饭了。
        贾母因问道：“跟着你娘吃了什 么好的了？” 宝玉笑道：“也没什么好的，我倒多吃了一碗饭。” 因问：“林姑娘 在那里？” 贾母道：“里头屋里呢。”}
}
{
    \eA{When Lin Tai-yü caught, with her ears, his protestations, and noticed with
        her eyes his state of mind, she unconsciously experienced an inward pang,
        and, much against her will, tears too besprinkled her cheeks; so, drooping
        her head, she kept silent. Her manner did not escape Pao-yü's notice. "I
        myself am aware," he speedily resumed, "that I'm worth nothing now;}
}
{
}

\Verse{20}
{
    \zA{宝玉进来，只见地下一个丫头吹熨斗，炕上 两个丫头打粉线，黛玉弯着腰拿剪子裁什么呢。 宝玉走进来，笑道：“哦!这是做 什么呢？
        才吃了饭，这么控着头，一会子又头疼了。” 黛玉并不理，只管裁他的。 有一个丫头说道：“那块绸子角儿还不好呢，再熨熨罢。” 黛玉便把剪子一撂，说
        道：“理他呢，过一会子就好了。” 宝玉听了，自是纳闷。 只见宝钗、探春等也来 了，和贾母说了一回话，宝钗也进来问：“妹妹做什么呢？”}
}
{
    \eA{but, however imperfect I may be, I could on no account presume to become
        guilty of any shortcoming with you cousin. Were I to ever commit the
        slightest fault, your task should be either to tender me advice and warn me
        not to do it again, or to blow me up a little, or give me a few whacks; and
        all this reproof I wouldn't take amiss.}
}
{
}

\Verse{21}
{
    \zA{因见林黛玉裁剪，笑 道：“越发能干了，连裁铰都会了。” 黛玉笑道：“这也不过是撒谎哄人罢了。”
        宝钗笑道：“我告诉你个笑话儿，才刚为那个药，我说了个不知道，宝兄弟心里就 不受用了。” 黛玉道：“理他呢，过会子就好了。” 宝玉向宝钗道：“老太太要抹
        骨牌，正没人，你抹骨牌去罢。” 宝钗听说，便笑道：“我是为抹骨牌才来么？” 说着便走了。 黛玉道：“你倒是去罢，这里有老虎，看吃了你！” 说着又裁。}
}
{
    \eA{But no one would have ever anticipated that you wouldn't bother your head in
        the least about me, and that you would be the means of driving me to my
        wits' ends, and so much out of my mind and off my head, as to be quite at a
        loss how to act for the best. In fact, were death to come upon me, I would
        be a spirit driven to my grave by grievances.}
}
{
}

\Verse{22}
{
    \zA{宝玉 见他不理，只得还陪笑说道：“你也去逛逛，再裁不迟。” 黛玉总不理。 宝玉便问 丫头们：“这是谁叫他裁的？”
        黛玉见问丫头们，便说道：“凭他谁叫我裁，也不 管二爷的事。” 宝玉方欲说话，只见有人进来，回说“外头有人请呢”。 宝玉听了， 忙撤身出来。
        黛玉向外头说道：“阿弥陀佛，赶你回来，我死了也罢了！” 宝玉来到外面，只见焙茗说：“冯大爷家请。” 宝玉听了，知道是昨日的话， 便说：“要衣裳去。”
        就自己往书房里来。}
}
{
    \eA{However much exalted bonzes and eminent Taoist priests might do penance,
        they wouldn't succeed in releasing my soul from suffering; for it would
        still be needful for you to clearly explain the facts, so that I might at
        last be able to come to life." After lending him a patient ear, Tai-yü
        suddenly banished from her memory all recollection of the occurrences of the
        previous night.}
}
{
}

\Verse{23}
{
    \zA{焙茗一直到了二门前等人，只见出来了 一个老婆子，焙茗上去说道：“宝二爷在书房里等出门的衣裳，你老人家进去带个 信儿。”
        那婆子啐道：“呸!放你娘的屁!宝玉如今在园里住着，跟他的人都在园里， 你又跑了这里来带信儿了！” 焙茗听了笑道：“骂的是，我也糊涂了！” 说着，一
        径往东边二门前来。 可巧门上小厮在甬路底下踢球，焙茗将原故说了，有个小厮跑 了进去，半日才抱了一个包袱出来，递给焙茗。}
}
{
    \eA{"Well, in that case," she said, "why did you not let a servant-girl open the
        door when I came over?" This question took Pao-yü by surprise. "What prompts
        you to say this?" he exclaimed. "If I have done anything of the kind, may I
        die at once." "Psha!" cried Tai-yü, "it's not right that you-should
        recklessly broach the subject of living or dying at this early morn! If you
        say yea, it's yea; and nay, it's nay;}
}
{
}

\Verse{24}
{
    \zA{回到书房里，宝玉换上，叫人备马， 只带着焙茗、锄药、双瑞、寿儿四个小厮去了。 一径到了冯紫英门口，有人报与冯紫英，出来迎接进去。 只见薛蟠早已在那里
        久候了，还有许多唱曲儿的小厮们，并唱小旦的蒋玉函，锦香院的妓女云儿。 大家 都见过了，然后吃茶。 宝玉擎茶笑道：“前儿说的‘幸与不幸’之事，我昼夜悬想，
        今日一闻呼唤即至。” 冯紫英笑道：“你们令姑表弟兄倒都心实。 前日不过是我的 设辞，诚心请你们喝一杯酒，恐怕推托，才说下这句话。}
}
{
    \eA{what use is there to utter such oaths!" "I didn't really see you come over,"
        protested Pao-yü. "Cousin Pao-ch'ai it was, who came and sat for a while and
        then left." After some reflection, Lin Tai-yü smiled. "Yes," she observed,
        "your servant-girls must, I fancy, have been too lazy to budge, grumpy and
        in a cross-grained mood; this is probable enough."}
}
{
}

\Verse{25}
{
    \zA{谁知都信了真了。” 说毕， 大家一笑。 然后摆上酒来，依次坐定。 冯紫英先叫唱曲儿的小厮过来递酒，然后叫 云儿也过来敬三钟。
        那薛蟠三杯落肚，不觉忘了情，拉着云儿的手笑道：“你把那 体己新鲜曲儿唱个我听，我喝一坛子，好不好？” 云儿听说，只得拿起琵琶来，唱 道：
        两个冤家，都难丢下，想着你来又惦记着他。 两个人形容俊俏都难描画，想昨 宵幽期私订在荼？ 架。 一个偷情，一个寻拿：拿住了三曹对案我也无回话。
        唱毕，笑道：“你喝一坛子罢了。”}
}
{
    \eA{"This is, I feel sure, the reason," answered Pao-yü, "so when I go back,
        I'll find out who it was, call them to task and put things right." "Those
        girls of yours;" continued Tai-yü, "should be given a lesson, but properly
        speaking it isn't for me to mention anything about it. Their present insult
        to me is a mere trifle;}
}
{
}

\Verse{26}
{
    \zA{薛蟠听说，笑道：“不值一坛，再唱好的来。” 宝玉笑道：“听我说罢：这么滥饮，易醉而无味。 我先喝一大海，发一个新令，
        有不遵者，连罚十大海，逐出席外，给人斟酒。” 冯紫英蒋玉函等都道：“有理， 有理。” 宝玉拿起海来，一气饮尽，说道：“如今要说‘悲’‘愁’‘喜’‘乐’
        四个字，却要说出‘女儿’来，还要注明这四个字的原故。 说完了，喝门杯，酒面 要唱一个新鲜曲子，酒底要席上生风一样东西――或古诗、旧对、《四书》《五经》
        成语。”}
}
{
    \eA{but were to-morrow some Miss Pao (precious) or some Miss Pei (jewel) or
        other to come, and were she to be subjected to insult, won't it be a grave
        matter?" While she taunted him, she pressed her lips, and laughed
        sarcastically. Pao-yü heard her remarks and felt both disposed to gnash his
        teeth with rage, and to treat them as a joke;}
}
{
}

\Verse{27}
{
    \zA{薛蟠不等说完，先站起来拦道：“我不来，别算我。 这竟是玩我呢！” 云 儿也站起来，推他坐下，笑道：“怕什么？ 这还亏你天天喝酒呢，难道连我也不及？
        我回来还说呢。 说是了罢，不是了不过罚上几杯，那里就醉死了你？ 如今一乱令， 倒喝十大海，下去斟酒不成？” 众人都拍手道：“妙！”
        薛蟠听说无法，只得坐了。 听宝玉说道：“女儿悲，青春已大守空闺。 女儿愁，悔教夫婿觅封侯。 女儿喜， 对镜晨妆颜色美。 女儿乐，秋千架上春衫薄。”}
}
{
    \eA{but in the midst of their colloquy, they perceived a waiting-maid approach
        and invite them to have their meal. Presently, the whole body of inmates
        crossed over to the front. "Miss," inquired Madame Wang at the sight of
        Tai-yü, "have you taken any of Dr. Pao's medicines? Do you feel any better?"
        "I simply feel so-so," replied Lin Tai-yü, "but grandmother Chia recommended
        me to go on taking Dr. Wang's medicines."}
}
{
}

\Verse{28}
{
    \zA{众人听了，都说道：“好！” 薛蟠 独扬着脸，摇头说：“不好，该罚。” 众人问：“如何该罚？” 薛蟠道：“他说的 我全不懂，怎么不该罚？”
        云儿便拧他一把，笑道：“你悄悄儿的想你的罢。 回来 说不出来，又该罚了。” 于是拿琵琶听宝玉唱道： 滴不尽相思血泪抛红豆，开不完春柳春花满画楼。
        睡不稳纱窗风雨黄昏后，忘 不了新愁与旧愁。 咽不下玉粒金波噎满喉，照不尽菱花镜里形容瘦。}
}
{
    \eA{"Mother," Pao-yü interposed, "you've no idea that cousin Lin's is an
        internal derangement; it's because she was born with a delicate physique
        that she can't stand the slightest cold. All she need do is to take a couple
        of closes of some decoction to dispel the chill; yet it's preferable that
        she should have medicine in pills."}
}
{
}

\Verse{29}
{
    \zA{展不开的眉头， 捱不明的更漏：呀!恰便似遮不住的青山隐隐，流不断的绿水悠悠。 唱完，大家齐声喝彩，独薛蟠说：“没板儿。” 宝玉饮了门杯，便拈起一片梨来，
        说道：“‘雨打梨花深闭门’。” 完了令。 下该冯紫英，说道：“女儿喜，头胎养了双生子。 女儿乐，私向花园掏蟋蟀。 女儿悲，儿夫染病在垂危。
        女儿愁，大风吹倒梳妆楼。” 说毕，端起酒来，唱道： 你是个可人，你是个多情，你是个刁钻古怪鬼灵精，你是个神仙也不灵。}
}
{
    \eA{"The other day," said Madame Wang, "the doctor mentioned the name of some
        pills, but I've forgotten what it is." "I know something about pills," put
        in Pao-yü; "he merely told her to take some pills or other called 'ginseng
        as-a-restorative-of-the-system.'" "That isn't it," Madame Wang demurred.
        "The 'Eight-precious-wholesome-to-mother' pills," Pao-yü proceeded, "or the
        'Left-angelica' or 'Right-angelica;'}
}
{
}

\Verse{30}
{
    \zA{我说 的话儿你全不信，只叫你去背地里细打听，才知道我疼你不疼! 唱完，饮了门杯，说道：“‘鸡声茅店月’。” 令完。
        下该云儿，云儿便说道：“女儿悲，将来终身倚靠谁？” 薛蟠笑道：“我的儿， 有你薛大爷在，你怕什么？” 众人都道：“别混他，别混他！” 云儿又道：“女儿
        愁，妈妈打骂何时休？” 薛蟠道：“前儿我见了你妈，还嘱咐他，不叫他打你呢。” 众人都道：“再多说的，罚酒十杯！”}
}
{
    \eA{if these also aren't the ones, they must be the 'Eight-flavour
        Rehmannia-glutinosa' pills." "None of these," rejoined Madame Wang, "for I
        remember well that there were the two words chin kang (guardians in
        Buddhistic temples)." "I've never before," observed Pao-yü, clapping his
        hands, "heard of the existence of chin kang pills;}
}
{
}

\Verse{31}
{
    \zA{薛蟠连忙自己打了一个嘴巴子，说道：“没 耳性，再不许说了。” 云儿又说：“女儿喜，情郎不舍还家里。 女儿乐，住了箫管 弄弦索。” 说完，便唱道：
        豆蔻花开三月三，一个虫儿往里钻。 钻了半日钻不进去，爬到花儿上打秋千。 肉儿小心肝，我不开了你怎么钻？ 唱毕，饮了门杯，说道：“‘桃之夭夭’。”
        令完，下该薛蟠。 薛蟠道：“我可要说了：女儿悲――”说了，半日不见说底下的。 冯紫英笑道： “悲什么？ 快说。”}
}
{
    \eA{but in the event of there being any chin kang pills, there must, for a
        certainty, be such a thing as P'u Sa (Buddha) powder." At this joke, every
        one in the whole room burst out laughing. Pao-ch'ai compressed her lips and
        gave a smile. "It must, I'm inclined to think," she suggested, "be the
        'lord-of-heaven-strengthen-the-heart' pills!"}
}
{
}

\Verse{32}
{
    \zA{薛蟠登时急的眼睛铃铛一般，便说道：“女儿悲――”又咳嗽 了两声，方说道：“女儿悲，嫁了个男人是乌龟。” 众人听了都大笑起来。 薛蟠道： “笑什么？
        难道我说的不是？ 一个女儿嫁了汉子，要做忘八，怎么不伤心呢？” 众人 笑的弯着腰说道：“你说的是!快说底下的罢。” 薛蟠瞪了瞪眼，又说道：“女儿
        愁――”说了这句，又不言语了。 众人道：“怎么愁？” 薛蟠道：“绣房钻出个大 马猴。” 众人哈哈笑道：“该罚，该罚!先还可恕，这句更不通了。”}
}
{
    \eA{"Yes, that's the name," Madame Wang laughed, "why, now, I too have become
        muddle-headed." "You're not muddle-headed, mother," said Pao-yü, "it's the
        mention of Chin kangs and Buddhas which confused you." "Stuff and nonsense!"
        ejaculated Madame Wang. "What you want again is your father to whip you!"
        "My father," Pao-yü laughed, "wouldn't whip me for a thing like this."}
}
{
}

\Verse{33}
{
    \zA{说着，便要 斟酒。 宝玉道：“押韵就好。” 薛蟠道：“令官都准了，你们闹什么！” 众人听说 方罢了。 云儿笑道：“下两句越发难说了，我替你说罢。”
        薛蟠道：“胡说!当真 我就没好的了？ 听我说罢：女儿喜，洞房花烛朝慵起。” 众人听了，都诧异道：“这 句何其太雅？” 薛蟠道：“女儿乐，一根？？
        往里戳。” 众人听了，都回头说道： “该死，该死!快唱了罢。” 薛蟠便唱道：“一个蚊子哼哼哼。” 众人都怔了，说 道：“这是个什么曲儿？”}
}
{
    \eA{"Well, this being their name," resumed Madame Wang, "you had better tell
        some one to-morrow to buy you a few." "All these drugs," expostulated
        Pao-yü, "are of no earthly use. Were you, mother, to give me three hundred
        and sixty taels, I'll concoct a supply of pills for my cousin, which I can
        certify will make her feel quite herself again before she has finished a
        single supply." "What trash!" cried Madame Wang.}
}
{
}

\Verse{34}
{
    \zA{薛蟠还唱道：“两个苍蝇嗡嗡嗡。” 众人都道：“罢， 罢，罢！” 薛蟠道：“爱听不听，这是新鲜曲儿，叫做‘哼哼韵’儿，你们要懒怠
        听，连酒底儿都免了，我就不唱。” 众人都道：“免了罢，倒别耽误了别人家。” 于是蒋玉函说道：“女儿悲，丈夫一去不回归。 女儿愁，无钱去打桂花油。 女
        儿喜，灯花并头结双蕊。 女儿乐，夫唱妇随真和合。” 说毕，唱道： 可喜你天生成百媚娇，恰便似活神仙离碧霄。 度青春，年正小； 配鸾凤，真也 巧。}
}
{
    \eA{"What kind of medicine is there so costly!" "It's a positive fact," smiled
        Pao-yü. "This prescription of mine is unlike all others. Besides, the very
        names of those drugs are quaint, and couldn't be enumerated in a moment;
        suffice it to mention the placenta of the first child; three hundred and
        sixty ginseng roots, shaped like human beings and studded with leaves; four
        fat tortoises; full-grown polygonum multiflorum;}
}
{
}

\Verse{35}
{
    \zA{呀!看天河正高，听谯楼鼓敲，剔银灯同入鸳帏悄。 唱毕，饮了门杯，笑道：“这诗词上我倒有限，幸而昨日见了一副对子，只记得这 句，可巧席上还有这件东西。”
        说毕，便干了酒，拿起一朵木樨来，念道：“‘花 气袭人知昼暖’。” 众人都倒依了完令，薛蟠又跳起来喧嚷道：“了不得，了不得，
        该罚，该罚!这席上并没有宝贝，你怎么说起宝贝来了？” 蒋玉函忙说道：“何曾 有宝贝？” 薛蟠道：“你还赖呢!你再说。” 蒋玉函只得又念了一遍。}
}
{
    \eA{the core of the Pachyma cocos, found on the roots of a fir tree of a
        thousand years old; and other such species of medicines. They're not, I
        admit, out-of-the-way things; but they are the most excellent among that
        whole crowd of medicines; and were I to begin to give you a list of them,
        why, they'd take you all quite aback.}
}
{
}

\Verse{36}
{
    \zA{薛蟠道：“这 ‘袭人’可不是宝贝是什么？ 你们不信只问他！” 说毕，指着宝玉。 宝玉没好意思 起来，说：“薛大哥，你该罚多少？”
        薛蟠道：“该罚，该罚！” 说着，拿起酒来， 一饮而尽。 冯紫英和蒋玉函等还问他原故，云儿便告诉了出来，蒋玉函忙起身陪罪。 众人都道：“不知者不作罪。”
        少刻，宝玉出席解手，蒋玉函随着出来，二人站在廊檐下，蒋玉函又赔不是。 宝玉见他妩媚温柔，心中十分留恋，便紧紧的攥着他的手，叫他：“闲了往我们那 里去。}
}
{
    \eA{The year before last, I at length let Hsüeh P'an have this recipe, after he
        had made ever so many entreaties during one or two years. When, however, he
        got the prescription, he had to search for another two or three years and to
        spend over and above a thousand taels before he succeeded in having it
        prepared. If you don't believe me, mother, you are at liberty to ask cousin
        Pao-ch'ai about it."}
}
{
}

\Verse{37}
{
    \zA{还有一句话问你，也是你们贵班中，有一个叫琪官儿的，他如今名驰天下， 可惜我独无缘一见。” 蒋玉函笑道：“就是我的小名儿。” 宝玉听说，不觉欣然跌
        足笑道：“有幸，有幸!果然名不虚传。 今儿初会，却怎么样呢？” 想了一想，向 袖中取出扇子，将一个玉？ 扇坠解下来，递给琪官，道：“微物不堪，略表今日之
        谊。” 琪官接了，笑道：“无功受禄，何以克当？ 也罢，我这里也得了一件奇物， 今日早起才系上，还是簇新，聊可表我一点亲热之意。”}
}
{
    \eA{At the mention of her name, Pao-ch'ai laughingly waved her hand. "I know
        nothing about it," she observed. "Nor have I heard anything about it, so
        don't tell your mother to ask me any questions." "Really," said Madame Wang
        smiling, "Pao-ch'ai is a good girl; she does not tell lies." Pao-yü was
        standing in the centre of the room. Upon hearing these words, he turned
        round sharply and clapped his hands.}
}
{
}

\Verse{38}
{
    \zA{说毕撩衣，将系小衣儿的 一条大红汗巾子解下来递给宝玉道：“这汗巾子是茜香国女国王所贡之物，夏天系 着肌肤生香，不生汗渍。 昨日北静王给的，今日才上身。
        若是别人，我断不肯相赠。 二爷请把自己系的解下来给我系着。” 宝玉听说，喜不自禁，连忙接了，将自己一 条松花汗巾解下来递给琪官。
        二人方束好，只听一声大叫：“我可拿住了！” 只见 薛蟠跳出来，拉着二人道：“放着酒不喝，两个人逃席出来，干什么？ 快拿出来我 瞧瞧。”}
}
{
    \eA{"What I stated just now," he explained, "was the truth; yet you maintain
        that it was all lies." As he defended himself, he casually looked round, and
        caught sight of Lin Tai-yü at the back of Pao-ch'ai laughing with tight-set
        lips, and applying her fingers to her face to put him to shame.}
}
{
}

\Verse{39}
{
    \zA{二人都道：“没有什么。” 薛蟠那里肯依，还是冯紫英出来才解开了。 复 又归坐饮酒，至晚方散。
        宝玉回至园中，宽衣吃茶，袭人见扇上的坠儿没了，便问他：“往那里去了？” 宝玉道：“马上丢了。” 袭人也不理论。 及睡时，见他腰里一条血点似的大红汗巾
        子，便猜着了八九分，因说道：“你有了好的系裤子了，把我的那条还我罢。” 宝 玉听说，方想起那汗巾子原是袭人的，不该给人。 心里后悔，口里说不出来，只得
        笑道：“我赔你一条罢。”}
}
{
    \eA{But Lady Feng, who had been in the inner rooms overseeing the servants
        laying the table, came out at once, as soon as she overheard the
        conversation. "Brother Pao tells no lies," she smilingly chimed in, "this is
        really a fact.}
}
{
}

\Verse{40}
{
    \zA{袭人听了，点头叹道：“我就知道你又干这些事了，也 不该拿我的东西给那些混帐人哪。 也难为你心里没个算计儿！” 还要说几句，又恐
        怄上他的酒来，少不得也睡了。 一宿无话。 次日天明方醒，只见宝玉笑道：“夜里失了盗也不知道，你瞧瞧裤子上。” 袭
        人低头一看，只见昨日宝玉系的那条汗巾子，系在自己腰里了，便知是宝玉夜里换 的，忙一顿就解下来，说道：“我不希罕这行子，趁早儿拿了去。” 宝玉见他如此，
        只得委婉解劝了一回。 袭人无法，暂且系上。}
}
{
    \eA{Some time ago cousin Hsüeh P'an came over in person and asked me for pearls,
        and when I inquired of him what he wanted them for, he explained that they
        were intended to compound some medicine with; adding, in an aggrieved way,
        that it would have been better hadn't he taken it in hand for he never had
        any idea that it would involve such a lot of trouble!}
}
{
}

\Verse{41}
{
    \zA{过后宝玉出去，终久解下来，扔在个 空箱子里了，自己又换了一条系着。 宝玉并未理论。 因问起：“昨日可有什么事情？” 袭人便回说：“二奶奶打发
        人叫了小红去了。 他原要等你来着，我想什么要紧，我就做了主，打发他去了。” 宝玉道：“很是。 我已经知道了，不必等我罢了。” 袭人又道：“昨儿贵妃打发夏
        太监出来送了一百二十两银子，叫在清虚观初一到初三打三天平安醮，唱戏献供， 叫珍大爷领着众位爷们跪香拜佛呢。 还有端午儿的节礼也赏了。”}
}
{
    \eA{When I questioned him what the medicine was, he returned for answer that it
        was a prescription of brother Pao's; and he mentioned ever so many
        ingredients, which I don't even remember. 'Under other circumstances,' he
        went on to say, 'I would have purchased a few pearls, but what are
        absolutely wanted are such pearls as have been worn on the head; and that's
        why I come to ask you, cousin, for some.}
}
{
}

\Verse{42}
{
    \zA{说着，命小丫头 来，将昨日的所赐之物取出来，却是上等宫扇两柄，红麝香珠二串，凤尾罗二端， 芙蓉簟一领。 宝玉见了，喜不自胜，问：“别人的也都是这个吗？”
        袭人道：“老 太太多着一个香玉如意，一个玛瑙枕。 老爷、太太、姨太太的，只多着一个香玉如 意。 你的和宝姑娘的一样。
        林姑娘和二姑娘、三姑娘、四姑娘只单有扇子和数珠儿， 别的都没有。 大奶奶、二奶奶他两个是每人两匹纱、两匹罗，两个香袋儿，两个锭 子药。”}
}
{
    \eA{If, cousin, you've got no broken ornaments at hand, in the shape of flowers,
        why, those that you have on your head will do as well; and by and bye I'll
        choose a few good ones and give them to you, to wear.' I had no other course
        therefore than to snap a couple of twigs from some flowers I have, made of
        pearls, and to let him take them away.}
}
{
}

\Verse{43}
{
    \zA{宝玉听了，笑道：“这是怎么个原故，怎么林姑娘的倒不和我的一样，倒是宝 姐姐的和我一样？ 别是传错了罢？” 袭人道：“昨儿拿出来，都是一分一分的写着
        签子，怎么会错了呢。 你的是在老太太屋里，我去拿了来了的。 老太太说了：明儿 叫你一个五更天进去谢恩呢。” 宝玉道：“自然要走一趟。”
        说着，便叫了紫鹃来： “拿了这个到你们姑娘那里去，就说是昨儿我得的，爱什么留下什么。” 紫鹃答应 了，拿了去。}
}
{
    \eA{One also requires a piece of deep red gauze, three feet in length of the
        best quality; and the pearls must be triturated to powder in a mortar."
        After each sentence expressed by lady Feng, Pao-yü muttered an invocation to
        Buddha. "The thing is as clear as sunlight now," he remarked. The moment
        lady Feng had done speaking, Pao-yü put in his word.}
}
{
}

\Verse{44}
{
    \zA{不一时回来，说：“姑娘说了，昨儿也得了，二爷留着罢。” 宝玉听 说，便命人收了。 刚洗了脸出来，要往贾母那里请安去，只见黛玉顶头来了，宝玉赶上去笑道：
        “我的东西叫你拣，你怎么不拣？” 黛玉昨日所恼宝玉的心事，早又丢开，只顾今 日的事了，因说道：“我没这么大福气禁受，比不得宝姑娘，什么‘金’哪‘玉’
        的，我们不过是个草木人儿罢了！”}
}
{
    \eA{"Mother," he added, "you should know that this is a mere makeshift, for
        really, according to the letter of the prescription, these pearls and
        precious stones should, properly speaking, consist of such as had been
        obtained from, some old grave and been worn as head-ornaments by some
        wealthy and honourable person of bygone days. But how could one go now on
        this account and dig up graves, and open tombs!}
}
{
}

\Verse{45}
{
    \zA{宝玉听他提出“金玉”二字来，不觉心里疑猜， 便说道：“除了别人说什么金什么玉，我心里要有这个想头，天诛地灭，万世不得 人身！”
        黛玉听他这话，便知他心里动了疑了，忙又笑道：“好没意思，白白的起 什么誓呢？ 谁管你什么金什么玉的！” 宝玉道：“我心里的事也难对你说，日后自 然明白。
        除了老太太、老爷、太太这三个人，第四个就是妹妹了。 有第五个人，我 也起个誓。” 黛玉道：“你也不用起誓，我很知道你心里有‘妹妹’。}
}
{
    \eA{Hence it is that such as are simply in use among living persons can equally
        well be substituted." "O-mi-to-fu!" exclaimed Madame Wang, after listening
        to him throughout. "That will never do, and what an arduous job to uselessly
        saddle one's self with;}
}
{
}

\Verse{46}
{
    \zA{但只是见了 ‘姐姐’，就把‘妹妹’忘了。” 宝玉道：“那是你多心，我再不是这么样的。” 黛玉道：“昨儿宝丫头他不替你圆谎，你为什么问着我呢？
        那要是我，你又不知怎 么样了！” 正说着，只见宝钗从那边来了，二人便走开了。 宝钗分明看见，只装没看见，低头过去了。 到了王夫人那里，坐了一回，然后
        到了贾母这边，只见宝玉也在这里呢。 宝钗因往日母亲对王夫人曾提过“金锁是个 和尚给的，等日后有玉的方可结为婚姻”等语，所以总远着宝玉。}
}
{
    \eA{for even though there be interred in some graves people, who've been dead
        for several hundreds of years, it wouldn't be a propitious thing were their
        corpses turned topsy-turvey now and the bones abstracted; just for the sake
        of preparing some medicine or other." Pao-yü thereupon addressed himself to
        Tai-yü. "Have you heard what was said or not?" he asked.}
}
{
}

\Verse{47}
{
    \zA{昨日见元春所赐 的东西，独他和宝玉一样，心里越发没意思起来。 幸亏宝玉被一个黛玉缠绵住了， 心心念念只惦记着黛玉，并不理论这事。
        此刻忽见宝玉笑道：“宝姐姐，我瞧瞧你 的那香串子呢？” 可巧宝钗左腕上笼着一串，见宝玉问他，少不得褪了下来。
        宝钗原生的肌肤丰泽，一时褪不下来，宝玉在傍边看着雪白的胳膊，不觉动了 羡慕之心。 暗暗想道：“这个膀子若长在林姑娘身上，或者还得摸一摸； 偏长在他
        身上，正是恨我没福。”}
}
{
    \eA{"And is there, pray, any likelihood that cousin Secunda would also follow in
        my lead and tell lies?" While saying this, his eyes were, albeit his face
        was turned towards Lin Tai-yü, fixed upon Pao-ch'ai. Lin Tai-yü pulled
        Madame Wang. "You just listen to him, aunt," she observed. "All because
        cousin Pao-ch'ai would not accommodate him by lying, he appeals to me."}
}
{
}

\Verse{48}
{
    \zA{忽然想起“金玉”一事来，再看看宝钗形容，只见脸若银 盆，眼同水杏，唇不点而含丹，眉不画而横翠，比黛玉另具一种妩媚风流，不觉又 呆了。
        宝钗褪下串子来给他，他也忘了接。 宝钗见他呆呆的，自己倒不好意思的， 起来扔下串子。 回身才要走，只见黛玉蹬着门槛子，嘴里咬着绢子笑呢。 宝钗道：
        “你又禁不得风吹，怎么又站在那风口里？” 黛玉笑道：“何曾不是在房里来着。 只因听见天上一声叫，出来瞧了瞧，原来是个呆雁。” 宝钗道：“呆雁在那里呢？}
}
{
    \eA{"Pao-yü has a great knack," Madame Wang said, "of dealing contemptuously
        with you, his cousin." "Mother," Pao-yü smilingly protested, "you are not
        aware how the case stands. When cousin Pao-ch'ai lived at home, she knew
        nothing whatever about my elder cousin Hsüeh P'an's affairs, and how much
        less now that she has taken up her quarters inside the garden? She, of
        course, knows less than ever about them!}
}
{
}

\Verse{49}
{
    \zA{我也瞧瞧。” 黛玉道：“我才出来，他就‘忒儿’的一声飞了。” 口里说着，将手 里的绢子一甩，向宝玉脸上甩来，宝玉不知，正打在眼上，“嗳哟”了一声。
        要知端的，下回分解。 (本章完)}
}
{
    \eA{Yet, cousin Lin just now stealthily treated my statements as lies, and put
        me to the blush." These words were still on his lips, when they perceived a
        waiting-maid, from dowager lady Chia's apartments, come in quest of Pao-yü
        and Lin Tai-yü to go and have their meal. Lin Tai-yü, however, did not even
        call Pao-yü, but forthwith rising to her feet, she went along, dragging the
        waiting-maid by the hand.}
}
{
}

\Verse{50}
{
    \zA{}
}
{
    \eA{"Let's wait for master Secundus, Mr. Pao, to go along with us," demurred the
        girl. "He doesn't want anything to eat," Lin Tai-yü replied; "he won't come
        with us, so I'll go ahead." So saying she promptly left the room. "I'll have
        my repast with my mother to-day," Pao-yü said. "Not at all," Madame Wang
        remarked, "not at all. I'm going to fast to-day, so it's only right and
        proper that you should go and have your own." "I'll also fast with you
        then," Pao-yü retorted. As he spoke, he called out to the servant to go
        back, and rushing up to the table, he took a seat. Madame Wang faced
        Pao-ch'ai and her companions. "You, girls," she observed, "had better have
        your meal, and let him have his own way!" "It's only right that you should
        go," Pao-ch'ai smiled.}
}
{
}

\Verse{51}
{
    \zA{}
}
{
    \eA{"Whether you have anything to eat or not, you should go over for a while to
        keep company to cousin Lin, as she will be quite distressed and out of
        spirits." "Who cares about her!" Pao-yü rejoined, "she'll get all right
        again after a time." Shortly, they finished their repast. But Pao-yü
        apprehended, in the first place, that his grandmother Chia, would be
        solicitous on his account, and longed, in the second, to be with Lin Tai-yü,
        so he hurriedly asked for some tea to rinse his mouth with. "Cousin
        Secundus," T'an Ch'un and Hsi Ch'un interposed with an ironic laugh, "what's
        the use of the hurry-scurry you're in the whole day long!}
}
{
}

\Verse{52}
{
    \zA{}
}
{
    \eA{Even when you're having your meals, or your tea, you're in this sort of
        fussy helter-skelter!" "Make him hurry up and have his tea," Pao-ch'ai
        chimed in smiling, "so that he may go and look up his cousin Lin. He'll be
        up to all kinds of mischief if you keep him here!" Pao-yü drank his tea.
        Then hastily leaving the apartment, he proceeded straightway towards the
        eastern court. As luck would have it, the moment he got near lady Feng's
        court, he descried lady Feng standing at the gateway. While standing on the
        step, and picking her teeth with an ear-cleaner, she superintended about ten
        young servant-boys removing the flower-pots from place to place. As soon as
        she caught sight of Pao-yü approaching, she put on a smiling face. "You come
        quite opportunely," she said;}
}
{
}

\Verse{53}
{
    \zA{}
}
{
    \eA{"walk in, walk in, and write a few characters for me." Pao-yü had no option
        but to follow her in. When they reached the interior of her rooms, lady Feng
        gave orders to a servant to fetch a pen, inkslab and paper. "Forty rolls of
        deep red ornamented satin," she began, addressing herself to Pao-yü, "forty
        rolls of satin with dragons; a hundred rolls of gauzes of every colour, of
        the finest quality; four gold necklaces…." "What's this?" Pao-yü shouted,
        "it is neither a bill; nor is it a list of presents, and in what style shall
        I write it?" Lady Feng remonstrated with him. "Just you go on writing," she
        said, "for, in fact, as long as I can make out what it means, it's all that
        is needed." Pao-yü at this response felt constrained to proceed with the
        writing. This our lady Feng put the paper by.}
}
{
}

\Verse{54}
{
    \zA{}
}
{
    \eA{As she did so, "I've still something more to tell you," she smilingly
        pursued, "but I wonder whether you will accede to it or not. There is in
        your rooms a servant-maid, Hsiao Hung by name, whom I would like to bring
        over into my service, and I'll select several girls to-morrow to wait on
        you; will this do?" "The servants in my quarters," answered Pao-yü, "muster
        a large crowd, so that, cousin, you are at perfect liberty to send for any
        one of them, who might take your fancy; what's the need therefore of asking
        me about it?" "If that be so," continued lady Feng laughingly, "I'll tell
        some one at once to go and bring her over." "Yes, she can go and fetch her,"
        acquiesced Pao-yü.}
}
{
}

\Verse{55}
{
    \zA{}
}
{
    \eA{While replying, he made an attempt to take his leave. "Come back," shouted
        lady Feng, "I've got something more to tell you." "Our venerable senior has
        sent for me," Pao-yü rejoined; "if you have anything to tell me you must
        wait till my return." After this explanation, he there and then came over to
        his grandmother Chia's on this side, where he found that they had already
        got through their meal. "Have you had anything nice to eat with your
        mother?" old lady Chia asked. "There was really nothing nice," Pao-yü
        smiled. "Yet I managed to have a bowl of rice more than usual." "Where's
        cousin Lin?" he then inquired. "She's in the inner rooms," answered his
        grandmother. Pao-yü stepped in.}
}
{
}

\Verse{56}
{
    \zA{}
}
{
    \eA{He caught sight of a waiting-maid, standing below, blowing into an iron, and
        two servant-girls seated on the stove-couch making a chalk line. Tai-yü with
        stooping head was cutting out something or other with a pair of scissors she
        held in her hand. Pao-yü advanced further in. "O! what's this that you are
        up to!" he smiled. "You have just had your rice and do you bob your head
        down in this way! Why, in a short while you'll be having a headache again!"
        Tai-yü, however, did not heed him in the least, but busied herself cutting
        out what she had to do. "The corner of that piece of satin is not yet
        right," a servant-girl put in. "You had better iron it again!" Tai-yü threw
        down the scissors. "Why worry yourself about it?" she said; "it will get
        quite right after a time."}
}
{
}

\Verse{57}
{
    \zA{}
}
{
    \eA{But while Pao-yü was listening to what was being said, and was inwardly
        feeling in low spirits, he became aware that Pao-ch'ai, T'an Ch'un and the
        other girls had also arrived. After a short chat with dowager lady Chia,
        Pao-ch'ai likewise entered the apartment to find out what her cousin Lin was
        up to. The moment she espied Lin Tai-yü engaged in cutting out something:
        "You have," she cried, "attained more skill than ever; for there you can
        even cut out clothes!" "This too," laughed Tai-yü sarcastically, "is a mere
        falsehood, to hoodwink people with, nothing more."}
}
{
}

\Verse{58}
{
    \zA{}
}
{
    \eA{"I'll tell you a joke," replied Pao-ch'ai smiling, "when I just now said
        that I did not know anything about that medicine, cousin Pao-yü felt
        displeased." "Who cares!" shouted Lin Tai-yü. "He'll get all right shortly."
        "Our worthy grandmother wishes to play at dominoes," Pao-yü thereupon
        interposed directing his remarks to Pao-ch'ai; "and there's no one there at
        present to have a game with her; so you'd better go and play with her."
        "Have I come over now to play dominoes!" promptly smiled Pao-ch'ai when she
        heard his suggestion. With this remark, she nevertheless at once quitted the
        room. "It would be well for you to go," urged Lin Tai-yü, "for there's a
        tiger in here; and, look out, he might eat you up." As she spoke, she went
        on with her cutting.}
}
{
}

\Verse{59}
{
    \zA{}
}
{
    \eA{Pao-yü perceived how loath she was to give him any of her attention, and he
        had no alternative but to force a smile and to observe: "You should also go
        for a stroll! It will be time enough by and bye to continue your cutting."
        But Tai-yü would pay no heed whatever to him. Pao-yü addressed himself
        therefore to the servant-girls. "Who has taught her how to cut out these
        things?" he asked. "What does it matter who taught me how to cut?" Tai-yü
        vehemently exclaimed, when she realised that he was speaking to the maids.
        "It's no business of yours, Mr. Secundus." Pao-yü was then about to say
        something in his defence when he saw a servant come in and report that there
        was some one outside who wished to see him. At this announcement, Pao-yü
        betook himself with alacrity out of the room. "O-mi-to-fu!"}
}
{
}

\Verse{60}
{
    \zA{}
}
{
    \eA{observed Tai-yü, turning outwards, "it wouldn't matter to you if you found
        me dead on your return!" On his arrival outside, Pao-yü discovered Pei Ming.
        "You are invited," he said, "to go to Mr. Feng's house." Upon hearing this
        message, Pao-yü knew well enough that it was about the project mooted the
        previous day, and accordingly he told him to go and ask for his clothes,
        while he himself wended his steps into the library. Pei Ming came forthwith
        to the second gate and waited for some one to appear. Seeing an old woman
        walk out, Pei Ming went up to her.}
}
{
}

\Verse{61}
{
    \zA{}
}
{
    \eA{"Our Master Secundus, Mr. Pao," he told her, "is in the study waiting for
        his out-door clothes; so do go in, worthy dame, and deliver the message."
        "It would be better," replied the old woman, "if you did not echo your
        mother's absurdities! Our Master Secundus, Mr. Pao, now lives in the garden,
        and all the servants, who attend on him, stay in the garden; and do you
        again come and bring the message here?" At these words, Pei Ming smiled.
        "You're quite right," he rejoined, "in reproving me, for I've become quite
        idiotic." So saying, he repaired with quick step to the second gate on the
        east side, where, by a lucky hit, the young servant-boys on duty, were
        kicking marbles on the raised road. Pei Ming explained to them the object of
        his coming. A young boy thereupon ran in.}
}
{
}

\Verse{62}
{
    \zA{}
}
{
    \eA{After a long interval, he, at length, made his appearance, holding, enfolded
        in his arms, a bundle of clothes, which he handed to Pei Ming, who then
        returned to the library. Pao-yü effected a change in his costume, and giving
        directions to saddle his horse, he only took along with him the four
        servant-boys, Pei Ming, Chu Lo, Shuang Jui and Shou Erh, and started on his
        way. He reached Feng Tzu-ying's doorway by a short cut. A servant announced
        his arrival, and Feng Tzu-ying came out and ushered him in. Here he
        discovered Hsüeh P'an, who had already been waiting a long time, and several
        singing-boys besides; as well as Chiang Yü-han, who played female roles, and
        Yün Erh, a courtesan in the Chin Hsiang court. The whole company exchanged
        salutations. They next had tea.}
}
{
}

\Verse{63}
{
    \zA{}
}
{
    \eA{"What you said the other day," smiled Pao-yü, raising his cup, "about good
        fortune coming out of evil fortune has preyed so much upon my mind, both by
        day and night, that the moment I received your summons I hurried to come
        immediately." "My worthy cousins," rejoined Feng Tzu-ying smiling. "You're
        all far too credulous! It's a mere hoax that I made use of the other day.
        For so much did I fear that you would be sure to refuse if I openly asked
        you to a drinking bout, that I thought it fit to say what I did. But your
        attendance to-day, so soon after my invitation, makes it clear, little
        though one would have thought it, that you've all taken it as pure gospel
        truth." This admission evoked laughter from the whole company.}
}
{
}

\Verse{64}
{
    \zA{}
}
{
    \eA{The wines were afterwards placed on the table, and they took the seats
        consistent with their grades. Feng Tzu-ying first and foremost called the
        singing-boys and offered them a drink. Next he told Yün Erh to also approach
        and have a cup of wine. By the time, however, that Hsüeh P'an had had his
        third cup, he of a sudden lost control over his feelings, and clasping Yün
        Erh's hand in his: "Do sing me," he smiled, "that novel ballad of your own
        composition; and I'll drink a whole jar full. Eh, will you?" This appeal
        compelled Yün Erh to take up the guitar. She then sang:  Lovers have I two.
        To set aside either I cannot bear.}
}
{
}

\Verse{65}
{
    \zA{}
}
{
    \eA{When my heart longs for thee to come,  It also yearns for him. Both are in
        form handsome and fair. Their beauty to describe it would be hard. Just
        think, last night, when at a silent hour, we met in secret, by  the trellis
        frame laden with roses white,  One to his feelings stealthily was giving
        vent,  When lo, the other caught us in the act,  And laying hands on us;
        there we three stood like litigants before  the  bar. And I had, verily, no
        word in answer for myself to give. At the close of her song, she laughed.
        "Well now," she cried, "down with that whole jar!" "Why, it isn't worth a
        jarful," smiled Hsüeh P'an at these words. "Favour us with some other good
        song!" "Listen to what I have to suggest," Pao-yü interposed, a smile on his
        lips.}
}
{
}

\Verse{66}
{
    \zA{}
}
{
    \eA{"If you go on drinking in this reckless manner, we will easily get drunk and
        there will be no fun in it. I'll take the lead and swallow a large cupful
        and put in force a new penalty; and any one of you who doesn't comply with
        it, will be mulcted in ten large cupfuls, in quick succession!" Speedily
        rising from the banquet, he poured the wine for the company. Feng Tzu-ying
        and the rest meanwhile exclaimed with one voice: "Quite right! quite right!"
        Pao-yü then lifted a large cup and drained it with one draught. "We will
        now," he proposed, "dilate on the four characters, 'sad, wounded, glad and
        joyful.' But while discoursing about young ladies, we'll have to illustrate
        the four states as well.}
}
{
}

\Verse{67}
{
    \zA{}
}
{
    \eA{At the end of this recitation, we'll have to drink the 'door cup' over the
        wine, to sing an original and seasonable ballad, while over the heel taps,
        to make allusion to some object on the table, and devise something with some
        old poetical lines or ancient scrolls, from the Four Books or the Five
        Classics, or with some set phrases." Hsüeh P'an gave him no time to finish.
        He was the first to stand up and prevent him from proceeding. "I won't join
        you, so don't count me; this is, in fact, done in order to play tricks upon
        me." Yün Erh, however, also rose to her feet and shoved him down into his
        seat. "What are you in such a funk for?" she laughed. "You're fortunate
        enough to be able to drink wine daily, and can't you, forsooth, even come up
        to me?}
}
{
}

\Verse{68}
{
    \zA{}
}
{
    \eA{Yet I mean to recite, by and bye, my own share. If you say what's right,
        well and good; if you don't, you will simply have to swallow several cups of
        wine as a forfeit, and is it likely you'll die from drunkenness? Are you,
        pray, going now to disregard this rule and to drink, instead, ten large
        cups; besides going down to pour the wine?" One and all clapped in applause.
        "Well said!" they shouted. After this, Hüeh P'an had no way out of it and
        felt compelled to resume his seat. They then heard Pao-yü recite:  A girl is
        sad,  When her spring-time of life is far advanced and she still occupies a
        vacant inner-room. A girl feels wounded in her heart,  When she regrets
        having allowed her better half to go abroad and win  a  marquisdom.}
}
{
}

\Verse{69}
{
    \zA{}
}
{
    \eA{A girl is glad,  When looking in the mirror, at the time of her morning
        toilette, she  finds her colour fair. A girl is joyful,  What time she sits
        on the frame of a gallows-swing, clad in a thin  spring gown. Having
        listened to him, "Capital!" one and all cried out in a chorus. Hsüeh P'an
        alone raised his face, shook his head and remarked: "It isn't good, he must
        be fined." "Why should he be fined?" demurred the party. "Because," retorted
        Hsüeh P'an, "what he says is entirely unintelligible to me. So how can he
        not be fined?" Yün Erh gave him a pinch.—"Just you quietly think of yours,"
        she laughed; "for if by and bye you are not ready you'll also have to bear a
        fine." In due course Pao-yü took up the guitar.}
}
{
}

\Verse{70}
{
    \zA{}
}
{
    \eA{He was heard to sing:  "When mutual thoughts arise, tears, blood-stained,
        endless drop, like  lentiles sown broadcast. In spring, in ceaseless bloom
        nourish willows and flowers around the  painted tower. Inside the
        gauze-lattice peaceful sleep flies, when, after dark, come  wind and rain.
        Both new-born sorrows and long-standing griefs cannot from memory  ever
        die! E'en jade-fine rice, and gold-like drinks they make hard to go down;
        they choke the throat. The lass has not the heart to desist gazing in the
        glass at her wan  face. Nothing can from that knitted brow of hers those
        frowns dispel; For hard she finds it patient to abide till the clepsydra
        will have  run its course. Alas! how fitly like the faint outline of a green
        hill which nought  can screen;}
}
{
}

\Verse{71}
{
    \zA{}
}
{
    \eA{Or like a green-tinged stream, which ever ceaseless floweth onward  far  and
        wide!" When the song drew to an end, his companions with one voice cried
        out: "Excellent!" Hsüeh P'an was the only one to find fault. "There's no
        metre in them," he said. Pao-yü quaffed the "opening cup," then seizing a
        pear, he added: "While the rain strikes the pear-blossom I firmly close the
        door," and thus accomplished the requirements of the rule. Feng Tzu-ying's
        turn came next. "A maid is glad." he commenced:  When at her first
        confinement she gives birth to twins, both sons. A maid is joyful,  When on
        the sly she to the garden creeps crickets to catch.}
}
{
}

\Verse{72}
{
    \zA{}
}
{
    \eA{A maid is sad,  When her husband some sickness gets and lies in a bad state.
        A maiden is wounded at heart,  When a fierce wind blows down the tower,
        where she makes her  toilette. Concluding this recitation, he raised the cup
        and sang:  "Thou art what one could aptly call a man. But thou'rt endowed
        with somewhat too much heart! How queer thou art, cross-grained and impish
        shrewd! A spirit too, thou couldst not be more shrewd. If all I say thou
        dost not think is true,  In secret just a minute search pursue; For then
        thou'lt know if I love thee or not." His song over, he drank the "opening
        cup" and then observed: "The cock crows when the moon's rays shine upon the
        thatchèd inn." After his observance of the rule followed Yün Erh's turn.}
}
{
}

\Verse{73}
{
    \zA{}
}
{
    \eA{A girl is sad, Yün Erh began,  When she tries to divine on whom she will
        depend towards the end of  life. "My dear child!" laughingly exclaimed Hsüeh
        P'an, "your worthy Mr. Hsüeh still lives, and why do you give way to fears?"
        "Don't confuse her!" remonstrated every one of the party, "don't muddle
        her!" "A maiden is wounded at heart." Yün Erh proceeded:  "When her mother
        beats and scolds her and never for an instant doth  desist." "It was only
        the other day," interposed Hsüeh P'an, "that I saw your mother and that I
        told her that I would not have her beat you." "If you still go on babbling,"
        put in the company with one consent, "you'll be fined ten cups." Hsüeh P'an
        promptly administered himself a slap on the mouth.}
}
{
}

\Verse{74}
{
    \zA{}
}
{
    \eA{"How you lack the faculty of hearing!" he exclaimed. "You are not to say a
        word more!" "A girl is glad," Yün Erh then resumed:  When her lover cannot
        brook to leave her and return home. A maiden is joyful,  When hushing the
        pan-pipe and double pipe, a stringed instrument she  thrums. At the end of
        her effusion, she at once began to sing:  "T'is the third day of the third
        moon, the nutmegs bloom; A maggot, lo, works hard to pierce into a flower;
        But though it ceaseless bores it cannot penetrate. So crouching on the buds,
        it swing-like rocks itself. My precious pet, my own dear little darling,  If
        I don't choose to open how can you steal in?" Finishing her song, she drank
        the "opening cup," after which she added: "the delicate peach-blossom," and
        thus complied with the exigencies of the rule.}
}
{
}

\Verse{75}
{
    \zA{}
}
{
    \eA{Next came Hsüeh P'an. "Is it for me to speak now?" Hsüeh P'an asked. "A
        maiden is sad…" But a long time elapsed after these words were uttered and
        yet nothing further was heard. "Sad for what?" Feng Tzu-ying laughingly
        asked. "Go on and tell us at once!" Hsüeh P'an was much perplexed. His eyes
        rolled about like a bell. "A girl is sad…" he hastily repeated. But here
        again he coughed twice before he proceeded. "A girl is sad." he said: "When
        she marries a spouse who is a libertine." This sentence so tickled the fancy
        of the company that they burst out into a loud fit of laughter. "What amuses
        you so?" shouted Hsüeh P'an, "is it likely that what I say is not correct?
        If a girl marries a man, who chooses to forget all virtue, how can she not
        feel sore at heart?"}
}
{
}

\Verse{76}
{
    \zA{}
}
{
    \eA{But so heartily did they all laugh that their bodies were bent in two. "What
        you say is quite right," they eagerly replied. "So proceed at once with the
        rest." Hsüeh P'an thereupon stared with vacant gaze. "A girl is grieved…."
        he added: But after these few words he once more could find nothing to say.
        "What is she grieved about?" they asked. "When a huge monkey finds its way
        into the inner room." Hsüeh P'an retorted. This reply set every one
        laughing. "He must be mulcted," they cried, "he must be mulcted. The first
        one could anyhow be overlooked; but this line is more unintelligible." As
        they said this, they were about to pour the wine, when Pao-yü smilingly
        interfered. "The rhyme is all right," he observed.}
}
{
}

\Verse{77}
{
    \zA{}
}
{
    \eA{"The master of the rules," Hsüeh P'an remarked, "approves it in every way,
        so what are you people fussing about?" Hearing this, the company eventually
        let the matter drop. "The two lines, that follow, are still more difficult,"
        suggested Yün Erh with a smile, "so you had better let me recite for you."
        "Fiddlesticks!" exclaimed Hsüeh P'an, "do you really fancy that I have no
        good ones! Just you listen to what I shall say. "A girl is glad,  When in
        the bridal room she lies, with flowery candles burning, and  she is loth to
        rise at morn." This sentiment filled one and all with amazement. "How
        supremely excellent this line is!" they ejaculated. "A girl is joyful,"
        Hsüeh P'an resumed, "During the consummation of wedlock."}
}
{
}

\Verse{78}
{
    \zA{}
}
{
    \eA{Upon catching this remark, the party turned their heads away, and shouted:
        "Dreadful! Dreadful! But quick sing your song and have done." Forthwith
        Hsüeh P'an sang: "A mosquito buzzes heng, heng, heng!" Every one was taken
        by surprise. "What kind of song is this?" they inquired. But Hsüeh P'an went
        on singing: "Two flies buzz weng, weng, weng." "Enough," shouted his
        companions, "that will do, that will do!" "Do you want to hear it or not?"
        asked Hsüeh P'an, "this is a new kind of song, called the 'Heng, heng air,'
        but if you people are not disposed to listen, let me off also from saying
        what I have to say over the heel-taps and I won't then sing." "We'll let you
        off! We'll let you off," answered one and all, "so don't be hindering
        others."}
}
{
}

\Verse{79}
{
    \zA{}
}
{
    \eA{"A maiden is sad," Chiang Yü-han at once began,  When her husband leaves
        home and never does return. A maiden is disconsolate,  When she has no money
        to go and buy some \_olea frangrans\_ oil. A maiden is glad,  When the wick
        of the lantern forms two heads like twin flowers on one  stem. A maiden is
        joyful,  When true conjugal peace prevails between her and her mate. His
        recital over, he went on to sing:  "How I love thee with those seductive
        charms of thine, heaven-born! In truth thou'rt like a living fairy from the
        azure skies! The spring of life we now enjoy; we are yet young in years. Our
        union is, indeed, a happy match! But. lo! the milky way doth at its zenith
        soar; Hark to the drums which beat around in the watch towers;}
}
{
}

\Verse{80}
{
    \zA{}
}
{
    \eA{So raise the silver lamp and let us soft under the nuptial curtain  steal."
        Finishing the song, he drank the "opening cup." "I know," he smiled, "few
        poetical quotations bearing on this sort of thing. By a stroke of good
        fortune, however, I yesterday conned a pair of antithetical scrolls; of
        these I can only remember just one line, but lucky enough for me the object
        it refers to figures as well on this festive board." This said he forthwith
        drained the wine, and, picking up a bud of a diminutive variety of \_olea
        fragrans\_, he recited:  "When the perfume of flowers wafts (hsi jen) itself
        into a man, he  knows the day is warm." The company unanimously conceded
        that the rule had been adhered to.}
}
{
}

\Verse{81}
{
    \zA{}
}
{
    \eA{But Hsüeh P'an once again jumped up. "It's awful, awful!" he bawled out
        boisterously; "he should be fined, he should be made to pay a forfeit;
        there's no precious article whatever on this table; how is it then that you
        introduce precious things?" "There was nothing about precious things!"
        Chiang Yü-han vehemently explained. "What I are you still prevaricating?"
        Hsüeh P'an cried, "Well, repeat it again!" Chiang Yü-han had no other course
        but to recite the line a second time. "Now is not Hsi Jen a precious thing?"
        Hsüeh P'an asked. "If she isn't, what is she? And if you don't believe me,
        you ask him about it," pointing, at the conclusion of this remark, at
        Pao-yü. Pao-yü felt very uncomfortable. Rising to his feet, "Cousin," he
        observed, "you should be fined heavily." "I should be! I should be!"}
}
{
}

\Verse{82}
{
    \zA{}
}
{
    \eA{Hsüeh P'an shouted, and saying this, he took up the wine and poured it down
        his throat with one gulp. Feng Tzu-ying, Chiang Yü-han and their companions
        thereupon asked him to explain the allusion. Yün Erh readily told them, and
        Chiang Yü-han hastily got up and pleaded guilty. "Ignorance," the party said
        with one consent, "does not amount to guilt." But presently Pao-yü quitted
        the banquet to go and satisfy a natural want and Chiang Yü-han followed him
        out. The two young fellows halted under the eaves of the verandah, and
        Chiang Yü-han then recommenced to make ample apologies. Pao-yü, however, was
        so attracted by his handsome and genial appearance, that he took quite a
        violent fancy to him; and squeezing his hand in a firm grip.}
}
{
}

\Verse{83}
{
    \zA{}
}
{
    \eA{"If you have nothing to do," he urged, "do let us go over to our place. I've
        got something more to ask you. It's this, there's in your worthy company
        some one called Ch'i Kuan, with a reputation extending at present throughout
        the world; but, unfortunately, I alone have not had the good luck of seeing
        him even once." "This is really," rejoined Chiang Yü-han with a smile, "my
        own infant name." This disclosure at once made Pao-yü quite exuberant, and
        stamping his feet he smiled. "How lucky! I'm in luck's way!" he exclaimed.
        "In very truth your reputation is no idle report. But to-day is our first
        meeting, and what shall I do?" After some thought, he produced a fan from
        his sleeve, and, unloosening one of the jade pendants, he handed it to Ch'i
        Kuan. "This is a mere trifle," he said.}
}
{
}

\Verse{84}
{
    \zA{}
}
{
    \eA{"It does not deserve your acceptance, yet it will be a small souvenir of our
        acquaintance to-day." Ch'i Kuan received it with a smile. "I do not
        deserve," he replied, "such a present. How am I worthy of such an honour!
        But never mind, I've also got about me here a strange thing, which I put on
        this morning; it is brand-new yet, and will, I hope, suffice to prove to you
        a little of the feeling of esteem which I entertain for you." With these
        protestations, he raised his garment, and, untying a deep red sash, with
        which his nether clothes were fastened, he presented it to Pao-yü. "This
        sash," he remarked, "is an article brought as tribute from the Queen of the
        Hsi Hsiang Kingdom.}
}
{
}

\Verse{85}
{
    \zA{}
}
{
    \eA{If you attach this round you in summer, your person will emit a fragrant
        perfume, and it will not perspire. It was given to me yesterday by the
        Prince of Pei Ching, and it is only to-day that I put it on. To any one
        else, I would certainly not be willing to present it. But, Mr. Secundus,
        please do unfasten the one you have on and give it to me to bind round me."
        This proposal extremely delighted Pao-yü. With precipitate haste, he
        accepted his gift, and, undoing the dark brown sash he wore, he surrendered
        it to Ch'i Kuan. But both had just had time to adjust their respective
        sashes when they heard a loud voice say: "Oh! I've caught you!" And they
        perceived Hsüeh P'an come out by leaps and bounds. Clutching the two young
        fellows, "What do you," he exclaimed, "leave your wine for and withdraw from
        the banquet.}
}
{
}

\Verse{86}
{
    \zA{}
}
{
    \eA{Be quick and produce those things, and let me see them!" "There's nothing to
        see!" rejoined the two young fellows with one voice. Hsüeh P'an, however,
        would by no means fall in with their views. And it was only Feng Tzu-ying,
        who made his appearance on the scene, who succeeded in dissuading him. So
        resuming their seats, they drank until dark, when the company broke up.
        Pao-yü, on his return into the garden, loosened his clothes, and had tea.
        But Hsi Jen noticed that the pendant had disappeared from his fan and she
        inquired of him what had become of it. "I must have lost it this very
        moment," Pao-yü replied. At bedtime, however, descrying a deep red sash,
        with spots like specks of blood, attached round his waist, Hsi Jen guessed
        more or less the truth of what must have transpired.}
}
{
}

\Verse{87}
{
    \zA{}
}
{
    \eA{"As you have such a nice sash to fasten your trousers with," Hsi Jen
        consequently said, "you'd better return that one of mine." This reminder
        made the fact dawn upon Pao-yü that the sash had originally been the
        property of Hsi Jen, and that he should by rights not have parted with it;
        but however much he felt his conscience smitten by remorse, he failed to see
        how he could very well disclose the truth to her. He could therefore only
        put on a smiling expression and add, "I'll give you another one instead."
        Hsi Jen was prompted by his rejoinder to nod her head and sigh. "I felt
        sure;" she observed; "that you'd go again and do these things! Yet you
        shouldn't take my belongings and bestow them on that low-bred sort of
        people. Can it be that no consideration finds a place in your heart?"}
}
{
}

\Verse{88}
{
    \zA{}
}
{
    \eA{She then felt disposed to tender him a few more words of admonition, but
        dreading, on the other hand, lest she should, by irritating him, bring the
        fumes of the wine to his head, she thought it best to also retire to bed.
        Nothing worth noticing occurred during that night. The next day, when she
        woke up at the break of day, she heard Pao-yü call out laughingly: "Robbers
        have been here in the night; are you not aware of it? Just you look at my
        trousers." Hsi Jen lowered her head and looked. She saw at a glance that the
        sash, which Pao-yü had worn the previous day, was bound round her own waist,
        and she at once realised that Pao-yü must have effected the change during
        the night; but promptly unbinding it, "I don't care for such things!" she
        cried, "quick, take it away!"}
}
{
}

\Verse{89}
{
    \zA{}
}
{
    \eA{At the sight of her manner, Pao-yü had to coax her with gentle terms. This
        so disarmed Hsi Jen, that she felt under the necessity of putting on the
        sash; but, subsequently when Pao-yü stepped out of the apartment, she at
        last pulled it off, and, throwing it away in an empty box, she found one of
        hers and fastened it round her waist. Pao-yü, however, did not in the least
        notice what she did, but inquired whether anything had happened the day
        before. "Lady Secunda," Hsi Jen explained, "dispatched some one and fetched
        Hsiao Hung away. Her wish was to have waited for your return; but as I
        thought that it was of no consequence, I took upon myself to decide, and
        sent her off." "That's all right!" rejoined Pao-yü.}
}
{
}

\Verse{90}
{
    \zA{}
}
{
    \eA{"I knew all about it, there was no need for her to wait." "Yesterday,"
        resumed Hsi Jen, "the Imperial Consort deputed the Eunuch Hsia to bring a
        hundred and twenty ounces of silver and to convey her commands that from the
        first to the third, there should be offered, in the Ch'ing Hsu temple,
        thanksgiving services to last for three days and that theatrical
        performances should be given, and oblations presented: and to tell our
        senior master, Mr. Chia Chen, to take all the gentlemen, and go and burn
        incense and worship Buddha.}
}
{
}

\Verse{91}
{
    \zA{}
}
{
    \eA{Besides this, she also sent presents for the dragon festival." Continuing,
        she bade a young servant-maid produce the presents, which had been received
        the previous day. Then he saw two palace fans of the best quality, two
        strings of musk-scented beads, two rolls of silk, as fine as the phoenix
        tail, and a superior mat worked with hibiscus. At the sight of these things,
        Pao-yü was filled with immeasurable pleasure, and he asked whether the
        articles brought to all the others were similar to his. "The only things in
        excess of yours that our venerable mistress has," Hsi Jen explained,
        "consist of a scented jade sceptre and a pillow made of agate.}
}
{
}

\Verse{92}
{
    \zA{}
}
{
    \eA{Those of your worthy father and mother, our master and mistress, and of your
        aunt exceed yours by a scented sceptre of jade. Yours are the same as Miss
        Pao's. Miss Lin's are like those of Misses Secunda, Tertia and Quarta, who
        received nothing beyond a fan and several pearls and none of all the other
        things. As for our senior lady, Mrs. Chia Chu, and lady Secunda, these two
        got each two rolls of gauze, two rolls of silk, two scented bags, and two
        sticks of medicine." After listening to her enumeration, "What's the reason
        of this?" he smiled. "How is it that Miss Lin's are not the same as mine,
        but that Miss Pao's instead are like my own? May not the message have been
        wrongly delivered?"}
}
{
}

\Verse{93}
{
    \zA{}
}
{
    \eA{"When they were brought out of the palace yesterday," Hsi Jen rejoined,
        "they were already divided in respective shares, and slips were also placed
        on them, so that how could any mistake have been made? Yours were among
        those for our dowager lady's apartments. When I went and fetched them, her
        venerable ladyship said that I should tell you to go there to-morrow at the
        fifth watch to return thanks. "Of course, it's my duty to go over," Pao-yü
        cried at these words, but forthwith calling Tzu Chüan: "Take these to your
        Miss Lin," he told her, "and say that I got them, yesterday, and that she is
        at liberty to keep out of them any that take her fancy." Tzu Chüan expressed
        her obedience and took the things away. After a short time she returned.}
}
{
}

\Verse{94}
{
    \zA{}
}
{
    \eA{"Miss Lin says," she explained, "that she also got some yesterday, and that
        you, Master Secundus, should keep yours." Hearing this reply, Pao-yü quickly
        directed a servant to put them away. But when he had washed his face and
        stepped out of doors, bent upon going to his grandmother's on the other
        side, in order to pay his obeisance, he caught sight of Lin Tai-yü coming
        along towards him, from the opposite direction. Pao-yü hurriedly walked up
        to her, "I told you," he smiled, "to select those you liked from my things;
        how is it you didn't choose any?" Lin Tai-yü had long before banished from
        her recollection the incident of the previous day, which had made her angry
        with Pao-yü, and was only exercised about the occurrence of this present
        occasion.}
}
{
}

\Verse{95}
{
    \zA{}
}
{
    \eA{"I'm not gifted with such extreme good fortune," she consequently answered,
        "as to be able to accept them. I can't compete with Miss Pao, in connection
        with whom something or other about gold or about jade is mentioned. We are
        simply beings connected with the vegetable kingdom." The allusion to the two
        words "gold and jade," aroused, of a sudden, much emotion in the heart of
        Pao-yü. "If beyond what people say about gold or jade," he protested, "the
        idea of any such things ever crosses my mind, may the heavens annihilate me,
        and may the earth extinguish me, and may I for ten thousand generations
        never assume human form!" These protestations convinced Lin Tai-yü that
        suspicion had been aroused in him. With all promptitude, she smiled and
        observed, "They're all to no use!}
}
{
}

\Verse{96}
{
    \zA{}
}
{
    \eA{Why utter such oaths, when there's no rhyme or reason! Who cares about any
        gold or any jade of yours!" "It would be difficult for me to tell you, to
        your face, all the secrets of my heart," Pao-yü resumed, "but by and bye
        you'll surely come to know all about them! After the three—my old
        grandmother, my father and my mother—you, my cousin, hold the fourth place;
        and, if there be a fifth, I'm ready to swear another oath." "You needn't
        swear any more," Lin Tai-yü replied, "I'm well aware that I, your younger
        cousin, have a place in your heart; but the thing is that at the sight of
        your elder cousin, you at once forget all about your younger cousin." "This
        comes again from over-suspicion!" ejaculated Pao; "for I'm not at all
        disposed that way."}
}
{
}

\Verse{97}
{
    \zA{}
}
{
    \eA{"Well," resumed Lin Tai-yü, "why did you yesterday appeal to me when that
        hussey Pao-ch'ai would not help you by telling a story? Had it been I, who
        had been guilty of any such thing, I don't know what you wouldn't have done
        again." But during their \_tête-a-tête\_, they espied Pao-ch'ai approach
        from the opposite direction, so readily they beat a retreat. Pao-ch'ai had
        distinctly caught sight of them, but pretending she had not seen them, she
        trudged on her way, with lowered head, and repaired into Madame Wang's
        apartments. After a short stay, she came to this side to pay dowager lady
        Chia a visit. With her she also found Pao-yü.}
}
{
}

\Verse{98}
{
    \zA{}
}
{
    \eA{Pao-ch'ai ever made it a point to hold Pao-yü aloof as her mother had in
        days gone by mentioned to Madame Wang and her other relatives that the gold
        locket had been the gift of a bonze, that she had to wait until such time as
        some suitor with jade turned up before she could be given in marriage, and
        other similar confidences. But on discovery the previous day that Yüan
        Ch'un's presents to her alone resembled those of Pao-yü, she began to feel
        all the more embarrassed. Luckily, however, Pao-yü was so entangled in Lin
        Tai-yü's meshes and so absorbed in heart and mind with fond thoughts of his
        Lin Tai-yü that he did not pay the least attention to this circumstance.}
}
{
}

\Verse{99}
{
    \zA{}
}
{
    \eA{But she unawares now heard Pao-yü remark with a smile: "Cousin Pao, let me
        see that string of scented beads of yours!" By a strange coincidence,
        Pao-ch'ai wore the string of beads round her left wrist so she had no
        alternative, when Pao-yü asked her for it, than to take it off. Pao-ch'ai,
        however, was naturally inclined to embonpoint, and it proved therefore no
        easy matter for her to get the beads off; and while Pao-yü stood by watching
        her snow-white arm, feelings of admiration were quickly stirred up in his
        heart. "Were this arm attached to Miss Lin's person," he secretly pondered,
        "I might, possibly have been able to caress it! But it is, as it happens,
        part and parcel of her body; how I really do deplore this lack of good
        fortune."}
}
{
}

\Verse{100}
{
    \zA{}
}
{
    \eA{Suddenly he bethought himself of the secret of gold and jade, and he again
        scanned Pao-ch'ai's appearance. At the sight of her countenance, resembling
        a silver bowl, her eyes limpid like water and almond-like in shape, her lips
        crimson, though not rouged, her eyebrows jet-black, though not pencilled,
        also of that fascination and grace which presented such a contrast to Lin
        Tai-yü's style of beauty, he could not refrain from falling into such a
        stupid reverie, that though Pao-ch'ai had got the string of beads off her
        wrist, and was handing them to him, he forgot all about them and made no
        effort to take them.}
}
{
}

\Verse{101}
{
    \zA{}
}
{
    \eA{Pao-ch'ai realised that he was plunged in abstraction, and conscious of the
        awkward position in which she was placed, she put down the string of beads,
        and turning round was on the point of betaking herself away, when she
        perceived Lin Tai-yü, standing on the door-step, laughing significantly
        while biting a handkerchief she held in her mouth. "You can't resist,"
        Pao-ch'ai said, "a single puff of wind; and why do you stand there and
        expose yourself to the very teeth of it?" "Wasn't I inside the room?"
        rejoined Lin Tai-yü, with a cynical smile. "But I came out to have a look as
        I heard a shriek in the heavens; it turned out, in fact, to be a stupid wild
        goose!" "A stupid wild goose!" repeated Pao-ch'ai. "Where is it, let me also
        see it!"}
}
{
}

\Verse{102}
{
    \zA{}
}
{
    \eA{"As soon as I got out," answered Lin Tai-yü, "it flew away with a 't'e-rh'
        sort of noise." While replying, she threw the handkerchief, she was holding,
        straight into Pao-yü's face. Pao-yü was quite taken by surprise. He was hit
        on the eye. "Ai-yah!" he exclaimed. But, reader, do you want to hear the
        sequel? In that case, listen to the circumstances, which will be disclosed
        in the next chapter.}
}
{
}
